\section{Self-Assessment and Development}
\label{sec:self_development}

\subsection{Skills Developed}

This project pushed me into unfamiliar domains. I learned MAVLink and ArduPilot~\cite{ardupilot,mavlink} from scratch---heartbeat negotiation, command acknowledgement, and autopilot configuration. Edge AI deployment (model export, TFLite~\cite{tflite} post-processing, Pi benchmarking) taught me that infrastructure around a model is where real complexity lies. Threading camera capture, inference, telemetry, and web streaming on a single-core Python process gave practical concurrency experience under GIL constraints. Field debugging on a Pi over SSH with no internet and a PuTTY terminal was perhaps the most underrated skill developed---problems that take minutes on a laptop took hours when the only feedback loop is \texttt{print} statements piped through a serial console.

\subsection{Time Management and Self-Management Failures}

I invested approximately 25--30 hours per week on this project across weeks 12--20, compared to an estimated 8--12 hours from most teammates. That imbalance was partly by necessity (I held the only programming expertise) and partly by choice (Section~\ref{sec:teamwork}). Work was structured around weekly milestones: weeks 12--15 on simulation, weeks 16--18 on Pi deployment, weeks 19--20 on field testing and retraining. A living ground-truth document (\texttt{CLAUDE.md}, 800+ lines) tracked every session's progress and generated the next session's task list. Version control discipline---feature branches, progressive commits, never breaking \texttt{main}---provided both a safety net and an audit trail.

\textbf{Where self-management failed.} I underestimated Pi integration by at least three weeks. My plan assumed ``same code runs everywhere''---and the architecture \textit{was} designed for that---but I did not deploy to the Pi until week~18. When I did, I discovered that Python~3.13 had broken \texttt{tflite-runtime} entirely (the package simply would not install), and that the IMX296 camera sensor outputs BGR despite the \texttt{picamera2} API labelling it RGB888. The TFLite fix required finding \texttt{ai-edge-litert}, an undocumented Google replacement package. The camera fix required testing all six channel permutations before identifying the correct one. Together, these consumed two full days---days I could not afford in week~18~\cite{koopman2019safe}.

The correction mechanism: I adopted a ``day-one deployment'' rule for all future hardware projects, and I applied it retroactively within this project. After the Pi debugging session, I immediately wrote a Docker-based Pi environment test (\texttt{Dockerfile.pi-test}) so that model loading could be verified on a Pi-like Linux image from my laptop. When weather cancelled our first flight day, I had pre-planned bench tests ready (servo checks, command sequences, camera stream verification), so the day remained productive rather than wasted.

\subsection{Strengths and Weaknesses: Two Sides of One Trait}

My core strength is \textbf{rapid prototyping and system integration}---decomposing problems into bounded modules and getting a working system fast. But this is inseparable from my primary weakness: \textbf{speed over inclusion}. The same instinct that produced a working state machine in two days also meant I built it alone, optimised for my own mental model, and presented teammates with a finished system rather than an opportunity to contribute. Edward observed that I ``built things faster than anyone could review them,'' which he meant as a compliment but which accurately describes the collaboration failure. The project manager noted that she sometimes felt ``out of the loop'' on technical decisions---not because I excluded her deliberately, but because I moved faster than I communicated. I suspect this pattern stems from my background: in Ukraine's engineering culture, individual technical competence is prized above collaborative process. Learning to value the slower, more inclusive approach as equally ``productive'' is an ongoing adjustment, not a one-time insight.

I communicate well through demonstrations and visual tools---the simulator and ground station conveyed the system architecture more effectively than any document---but I am weaker at spontaneous verbal explanation, especially when asked to slow down and simplify. This matters: a system only one person understands is a single point of failure, regardless of how well it works.

A second weakness is \textbf{delayed hardware testing}: I treated simulation as sufficient validation for too long. The architecture was genuinely cross-platform, but platform-specific bugs (Python~3.13, camera colour, serial baud rates) are invisible in simulation by definition. This is a pattern I recognise from previous projects---I gravitate toward the comfortable development environment and defer the uncomfortable one.

\subsection{Helping Teammates Develop}

Section~\ref{sec:teamwork} describes the collaboration structure; here I reflect on whether I developed people, not just output. There were genuine attempts: I walked the pilot through the ground station interface field-by-field so he understood what each telemetry value meant and could make independent abort decisions. I explained the state machine to the project manager using the simulator as a visual aid, which enabled her to write the system description section of the group report accurately without my editing. I paired with Edward on FOV calibration, explaining the trigonometry so he could run calibration independently on flight day.

But I also recognise the pattern Gibbs~\cite{gibbs1988} would flag: much of my ``feedback'' to teammates was corrective (rerouting wiring, adding missing risk items) rather than developmental. I reviewed the hardware lead's Cube wiring and pointed out a ground loop risk---useful, but it was me fixing his work, not teaching him to spot the issue himself. The difference matters for team capability, and I did not consistently achieve the latter.

\subsection{Development Programme}

Based on reflective analysis using Gibbs' cycle~\cite{gibbs1988} and Kolb's experiential learning model~\cite{kolb1984}, I identify three development actions, with evidence of implementation already begun within this project:

\begin{enumerate}
  \item \textbf{Collaborative architecture from day one.} Define contribution points and interface contracts \textit{before} implementation. I began applying this mid-project: \texttt{vision.py}'s single-function interface (\texttt{detect\_in\_image}) and numbered test scripts (0a--4) were deliberate attempts to lower the barrier. The lesson is that these must exist from week~1, not retrofitted in week~15.

  \item \textbf{Day-one hardware deployment.} Run a trivial script on target hardware within the first week to surface platform issues early. Already implemented: after the Python~3.13 incident, I created a Docker Pi test and added a preflight connectivity checker (\texttt{preflight.py}) that verifies camera, model, and Cube connection before any mission script runs.

  \item \textbf{Structured communication cadence.} Adopt a ``one sentence, one diagram, one demo'' framework for architecture explanations, and schedule regular check-ins rather than relying on teammates to ask. I partially implemented this through the simulator demo at PDR and the ground station's browser interface, but I did not establish a regular rhythm of proactive updates---teammates still had to ask for status.
\end{enumerate}

\subsection{Connection to Career Goals}

I came to Bristol from Ukraine to develop robotics skills I can bring back to help my country. This project directly involves technologies with humanitarian and defence applications---autonomous search, AI detection~\cite{yolov8}, and edge computing for field-deployable systems. The skills developed---MAVLink integration, edge AI, safety-critical design~\cite{leveson2011}---transfer to autonomous systems Ukraine needs for mine clearance, reconnaissance, and disaster response. The self-management lessons are equally transferable: in a conflict zone, delayed hardware testing is not a schedule inconvenience but a mission failure.

If I had to summarise this project's lesson: \textit{the quality of an autonomous system is determined not by its algorithms' cleverness, but by the rigour of its testing and the honesty of its failure analysis.}
